\subsection{\texorpdfstring{例1：复形 \(\mathbf{S}(\mathbf{L})\otimes_{\mathbf{A}}\boldsymbol{\Lambda}(\mathbf{L})\)}{例1：复形 \textbackslash mathbf\{S\}(\textbackslash mathbf\{L\})\textbackslash otimes\_\{\textbackslash mathbf\{A\}\}\textbackslash boldsymbol\{\textbackslash Lambda\}(\textbackslash mathbf\{L\})}}

设 A 为环，L 为 A-模， \(\mathbf{S}(\mathrm{L})\) 为其对称代数， \(\mathbf{S}(\mathrm{L}) \otimes_{\mathrm{A}} \mathrm{L}\) 为经纯量扩张得到的 \(\mathbf{S}(\mathrm{L})\) -模， \(u: \mathbf{S}(\mathrm{L}) \otimes_{\mathrm{A}} \mathrm{L} \to \mathbf{S}(\mathrm{L})\) 为线性形式，使得 \(u(s \otimes x) = sx\) 对 \(s \in \mathbf{S}(\mathrm{L})\) 、 \(x \in \mathrm{L}\) 成立。根据 \(\mathbf{S}(\mathrm{L})\) -模的典范同构（III，第 83 页，命题 8）

\[
\boldsymbol {\Lambda} _ {\mathbf {S} (L)} (\mathbf {S} (L) \otimes_ {A} L) \rightarrow \mathbf {S} (L) \otimes_ {A} \boldsymbol {\Lambda} (L),
\]

复形 \(\mathbf{K}^{\mathbf{S}(\mathrm{L})}(u)\) 的微分被搬运到映射

\[
d: \mathbf {S} (\mathrm{L}) \otimes_ {\mathrm{A}} \boldsymbol {\Lambda} (\mathrm{L}) \rightarrow \mathbf {S} (\mathrm{L}) \otimes_ {\mathrm{A}} \boldsymbol {\Lambda} (\mathrm{L})
\]

使得对 L 中的 \(x_{1},\ldots,x_{p},y_{1},\ldots,y_{q}\) ，我们有

\[
\begin{array}{l} d \left(\left(x _ {1} \dots x _ {p}\right) \otimes \left(y _ {1} \wedge \dots \wedge y _ {q}\right)\right) \\ = \sum_ {i = 1} ^ {q} (- 1) ^ {i + 1} y _ {i} x _ {1} \dots x _ {p} \otimes \left(y _ {1} \wedge \dots \wedge y _ {i - 1} \wedge y _ {i + 1} \wedge \dots \wedge y _ {q}\right). \end{array}
\]

注意 \(d\) 把 \(\mathbf{S}^{p}(\mathrm{L})\otimes\symbf{\Lambda}^{q}(\mathrm{L})\) 映到 \(\mathbf{S}^{p+1}(\mathrm{L})\otimes\symbf{\Lambda}^{q-1}(\mathrm{L})\) ，因此 A-模复形 \(\mathbf{S}(\mathrm{L})\otimes\symbf{\Lambda}(\mathrm{L})\) 分解为下述图表所描述的复形的直和：

\[
\left(\mathcal {E} _ {n}\right): 0 \rightarrow \mathbf {S} ^ {0} L \otimes_ {\mathrm{A}} \boldsymbol {\Lambda} ^ {n} L \rightarrow \mathbf {S} ^ {1} L \otimes_ {\mathrm{A}} \boldsymbol {\Lambda} ^ {n - 1} L \rightarrow \dots \rightarrow \mathbf {S} ^ {n} L \otimes_ {\mathrm{A}} \boldsymbol {\Lambda} ^ {0} L \rightarrow 0, n \in \mathbf {N}.
\]

若 A-模 L 是有限族 \((\mathbf{L}_{i})_{i\in I}\) 的直和，其中 I 是全序的，则典范双射

\[
\bigotimes_ {i \in I} ^ {\varepsilon} \left(\mathbf {S} (L _ {i}) \otimes_ {A} \boldsymbol {\Lambda} (L _ {i})\right) \to \mathbf {S} (L) \otimes_ {A} \boldsymbol {\Lambda} (L)
\]

是 A-模复形的同构（这由 X，第 148 页，命题 2 或上面的公式 (9) 推出）。

命题 3. --- 若 A-模 L 是平坦的，则上述序列 \((\mathcal{E}_n)\) 对 \(n > 0\) 是正合的。

\begin{enumerate}
\def\labelenumi{\alph{enumi})}
\tightlist
\item
  首先注意，若 \(p_{L}\) 是复合同态
\end{enumerate}

\[
\mathbf {S} (\mathrm{L}) \otimes \boldsymbol {\Lambda} (\mathrm{L}) \xrightarrow {\alpha} \mathbf {S} ^ {0} (\mathrm{L}) \otimes \boldsymbol {\Lambda} ^ {0} (\mathrm{L}) \xrightarrow {\beta} \mathrm{A},
\]

其中 \(\alpha\) 是诸典范投影的张量积， \(\beta\) 是典范同构；只需证明 \(\mathbf{H}(p_{\mathrm{L}})\) 是双射的。

\begin{enumerate}
\def\labelenumi{\alph{enumi})}
\setcounter{enumi}{1}
\item
  若L=0或L=A，命题显然成立。
\item
  假设 L 是有限秩的自由模；将它写成秩 1 的自由 A-模的直和 \(L_{1} \oplus \ldots \oplus L_{n}\) 。根据命题前的注记，复形 \(\mathbf{S}(\mathrm{L}) \otimes \symbf{\Lambda}(\mathrm{L})\) 同构于 n 个自由复形 \(\mathbf{S}(\mathrm{L}_{i}) \otimes \symbf{\Lambda}(\mathrm{L}_{i})\) 的张量积，后者的同调由 b) 知是自由的。由 X，第 79 页，推论 4，典范同态
\end{enumerate}

\[
\gamma : \bigoplus_ {i = 1} ^ {n} \mathrm{H} (\mathbf {S} (\mathrm{L} _ {i}) \otimes \boldsymbol {\Lambda} (\mathrm{L} _ {i})) \rightarrow \mathrm{H} (\mathbf {S} (\mathrm{L}) \otimes \boldsymbol {\Lambda} (\mathrm{L}))
\]

是双射的。由 b) 知 \(\mathrm{H}(p_{\mathrm{L}_i})\) 对所有 \(i\) 都是双射的。由于 \(\bigotimes_{i=1}^{n} \mathrm{H}(p_{\mathrm{L}_i}) = \mathrm{H}(p_{\mathrm{L}}) \circ \gamma\) ，

\[
\mathrm{H} (p _ {\mathrm{L}}) \text {   is   bijective }  .
\]

\begin{enumerate}
\def\labelenumi{\alph{enumi})}
\setcounter{enumi}{3}
\tightlist
\item
  在一般情况下，L 是由有限秩自由模构成的滤过归纳系 \((\mathbf{L}_{i})_{i\in I}\) 的归纳极限（X，第 14 页，定理 1）。由于典范双射同态
\end{enumerate}

\[
\varinjlim \mathbf {S} (\mathrm{L} _ {i}) \otimes \boldsymbol {\Lambda} (\mathrm{L} _ {i}) \rightarrow \mathbf {S} (\mathrm{L}) \otimes \boldsymbol {\Lambda} (\mathrm{L})
\]

是复形的同构，命题由 X，第 28 页，命题 1 得出。

注记。--- 1) 我们将在下文（X，第 158 页，例）看到上述 c) 部分的另一个证明。

\begin{enumerate}
\def\labelenumi{\arabic{enumi})}
\setcounter{enumi}{1}
\tightlist
\item
  如果 A 是一个 Q-代数，那么命题 3 的结论在不对 L 作任何假设的情况下仍然成立（参见 X，第 206 页，习题 1）。
\end{enumerate}

设 G 是一个群， \(\rho: G \to \mathbf{GL}(L)\) 是 G 在平坦 A-模 L 中的一个线性表示。则 \((\mathcal{E}_n)\) 是线性表示的正合列。假设 L 是有限生成投射模，并记 \(R_A(G)\) 为 G 在有限生成投射 A-模中的表示环。由命题 3 可知，我们在 \(R_{A}(G)\) 中有关系

\[
\sum_ {i = 0} ^ {n} (- 1) ^ {i} \left[ \mathbf {S} ^ {i} (\mathrm{L}) \right] \left[ \boldsymbol {\Lambda} ^ {n - i} (\mathrm{L}) \right] = 0, \quad n > 0.\tag{10}
\]

如果我们考虑形式幂级数

\[
s (\mathrm{T}) = \sum_ {i = 0} ^ {\infty} \left[ \mathbf {S} ^ {i} (\mathrm{L}) \right] \mathrm{T} ^ {i} \in \mathrm{R} _ {\mathbf {A}} (\mathrm{G}) [ [ \mathrm{T} ] ],
\]

\[
\lambda (\mathrm{T}) = \sum_ {i = 0} ^ {\infty} \left[ \boldsymbol {\Lambda} ^ {i} (\mathrm{L}) \right] \mathrm{T} ^ {i} \in \mathrm{R} _ {\mathbf {A}} (\mathrm{G}) [ [ \mathrm{T} ] ],
\]

关系式 (10) 可写作

\begin{enumerate}
\def\labelenumi{(\arabic{enumi})}
\setcounter{enumi}{10}
\tightlist
\item
\end{enumerate}

\[
s (\mathrm{T}) \lambda (- \mathrm{T}) = 1 _ {*}.
\]

